\documentclass[8pt,letterpaper]{article}

\usepackage[margin=0.30in,top=0.26in,bottom=0.26in]{geometry}
\usepackage[T1]{fontenc}
\usepackage{multicol}
\usepackage{enumitem}
\usepackage{booktabs}
\usepackage{array}
\usepackage{titlesec}

\pagestyle{empty}
\setlength{\parindent}{0pt}
\setlength{\parskip}{0pt}
\setlength{\columnsep}{8pt}
\setlength{\tabcolsep}{3pt}
\renewcommand{\arraystretch}{0.95}
\setlist[itemize]{leftmargin=9pt,itemsep=0pt,topsep=0pt,parsep=0pt}
\setlist[enumerate]{leftmargin=11pt,itemsep=0pt,topsep=0pt,parsep=0pt}
\titlespacing*{\section}{0pt}{0.45ex}{0.25ex}
\titlespacing*{\subsection}{0pt}{0.35ex}{0.15ex}

\newcommand{\term}[1]{\textbf{#1}}

\begin{document}
\fontsize{7.2pt}{8.1pt}\selectfont
\raggedbottom

\begin{center}
{\Large\bfseries Strategic Management Study Guide}\\[-1pt]
{UCM-integrated \textbullet{} Thursday 21st \textbullet{} printer-friendly compact version}\\[-1pt]
{Use this as an exam cram sheet: framework, evidence, implication, recommendation}
\end{center}

\begin{multicols}{2}

\section*{Exam pattern}
The fastest way to build a solid answer is to use one framework, apply it to the case, and finish with a recommendation. A good structure is: name the issue, choose the framework, give two or three evidence points, and close with the strategic implication. The exam rewards clear reasoning more than decorative formatting.

\term{Master pattern:} From a [framework] perspective, [force or concept] is high or low because [reason], which makes the industry attractive or unattractive for [who]; therefore [strategic implication].

\term{Stock openers:} This is about industry attractiveness. This is about internal advantage. The key issue is whether the strategy is coherent.

\term{How to score well:} Use the class vocabulary, name the framework explicitly, avoid generic statements, and always end with what the firm should do next.

\section*{1. Strategy fundamentals}
Strategy is the integrated set of actions a firm uses to exploit its competencies and create advantage. The strategic job is not just to decide what to do, but to make sure the choice fits the firm, the market, and the available resources.

There are three levels of strategy. \term{Corporate strategy} decides where the firm competes overall. \term{Business strategy} decides how a particular business unit competes. \term{Functional strategy} translates the plan into actions in marketing, operations, finance, and people management.

The planning hierarchy is simple: mission and vision set direction, strategic objectives define what the firm wants to achieve, strategy explains how it will get there, tactics turn that into specific actions, and operations handle day-to-day execution.

\term{Mission vs. vision:} Mission explains the current purpose and what the firm does now. Vision describes the future state or aspiration. Values are the principles that should guide behavior while the strategy is executed.

\section*{2. External analysis}
\term{PESTEL} is the macro scan: Political, Economic, Social, Technological, Environmental, and Legal factors. Use it to spot changes in the environment before they become obvious problems.

Think of PESTEL as the outside world changing the rules of the game. In an exam answer, macro trends matter when they alter demand, costs, regulation, or technology adoption.

\term{Porter's Five Forces} measure industry attractiveness. Strong forces reduce profit potential; weak forces make the industry easier to earn money in.

\begin{tabular}{@{}p{2.1cm}p{4.4cm}@{}}
\toprule
\term{Force} & \term{What to ask} \\
\midrule
New entrants & How easy is it for new firms to enter? \\
Suppliers & How much power do suppliers have over price or quality? \\
Buyers & Can customers force lower prices or better terms? \\
Substitutes & What alternative products can satisfy the same need? \\
Rivalry & How intense is competition among existing firms? \\
\bottomrule
\end{tabular}

When you explain the forces, push past naming them. Say why the force is strong, what evidence shows that, and what it means for margins, bargaining power, entry, or growth. Example: if buyer power is high, price pressure rises and the firm may need differentiation or switching costs.

\term{SWOT} should not become a list of random words. Strengths and weaknesses are internal; opportunities and threats are external. The useful part is the implication: each point should lead to a strategic action, not just a label.

Best practice is to convert SWOT into a TOWS-style action: use strengths to seize opportunities, use strengths to reduce threats, fix weaknesses that block opportunities, and defend against threats that exploit weaknesses.

\section*{3. Internal analysis and advantage}
\term{RBV} says that advantage comes from resources and capabilities that are valuable, hard to copy, and used well by the organization. \term{VRIO} is the clean test: Valuable, Rare, Inimitable, Organized. If a resource passes all four, it can support sustained competitive advantage.

The point of VRIO is not to recite letters. It is to ask whether a resource helps the firm exploit an opportunity or neutralize a threat, whether competitors also have it, whether it can be copied, and whether the firm is set up to capture the value.

\term{Value chain} analysis helps you find where value is actually created. Primary activities and support activities can explain cost leadership, differentiation, or a hidden bottleneck that competitors cannot easily match.

Primary activities are inbound logistics, operations, outbound logistics, marketing and sales, and service. Support activities are firm infrastructure, HR, technology development, and procurement. Use these to identify where the firm is strong or weak.

Competitive advantage usually takes one of three forms: \term{cost leadership}, \term{differentiation}, or \term{focus}. The warning is the same one from class: do not end up stuck in the middle, where the firm is neither the cheapest nor the most distinctive.

Cost leadership usually needs scale, process discipline, tight cost control, and efficient supply chains. Differentiation usually comes from brand, design, service, quality, technology, or customer experience. Focus means serving a narrow segment better than broad rivals can.

\term{Strategy Clock} is useful when you want to map price and perceived value together. The main idea is that firms can compete through low price, hybrid value, differentiation, or focused differentiation, depending on the target market and the offer.

If the firm is on the low-price end, the answer usually stresses efficiency and volume. If it is on the high-value end, the answer usually stresses brand, uniqueness, and customer willingness to pay.

\section*{4. Growth and corporate strategy}
\term{Ansoff} gives four growth paths: market penetration, product development, market development, and diversification. Risk rises as both the product and the market become new. Penetration is the safest move; diversification is the riskiest.

Use Ansoff to explain direction, not just categories. Market penetration means more sales to current customers or markets. Product development means new products for current customers. Market development means current products in new markets. Diversification means new products for new markets and is the most uncertain.

\term{Vertical integration} means moving backward toward suppliers or forward toward customers to control more of the value chain. It makes sense when supplier power, customer power, or transaction costs are a problem.

Backward integration helps when inputs are expensive or unreliable. Forward integration helps when the firm wants more control over distribution, pricing, or the customer relationship.

\term{Horizontal integration} is acquisition or merger with rivals or related firms to gain scale, capabilities, or market power.

This can reduce rivalry, create synergies, or widen the customer base, but it also creates integration risk, culture clashes, and possible regulatory scrutiny.

\term{Related diversification} uses shared capabilities across businesses. \term{Unrelated diversification} spreads risk but often weakens focus.

Related diversification is usually easier to justify in an exam because it creates synergies in brand, technology, channels, or capabilities. Unrelated diversification is harder to manage unless there is a strong financial reason.

\term{BCG matrix} is a portfolio tool for allocating resources. \term{Stars} deserve investment, \term{Cash Cows} generate cash and can be harvested, \term{Question Marks} need selective investment or exit decisions, and \term{Dogs} usually deserve divestment unless they serve a strategic purpose.

Use the BCG matrix to show where cash comes from and where it should go. The exam-friendly interpretation is: fund Stars and promising Question Marks with Cash Cow profits, and avoid wasting resources on weak Dogs unless they protect the broader portfolio.

\term{Business Model Canvas} is a reminder that strategy also has to work at the level of value proposition, customer segments, channels, relationships, revenue, key resources, key activities, partners, and cost structure.

In a case, use the canvas to ask: what is the value proposition, who is the customer, how does the firm deliver, how does it earn money, and what resources or partners make it possible?

\section*{5. Implementation and control}
The \term{Balanced Scorecard} translates strategy into execution through four lenses: financial, customer, internal process, and learning and growth. It is useful because it connects long-term strategy with measurable performance.

The main advantage of the scorecard is balance. It stops managers from focusing only on short-term financials and forces them to track customer satisfaction, internal efficiency, and capability building as well.

Strategies fail when structure, incentives, resources, and monitoring do not line up. A good idea with bad execution is still a bad strategy in practice.

Common failure points are unclear responsibilities, too many priorities, weak communication, poor coordination between functions, and no way to track progress.

The control loop is straightforward: set targets, measure performance, compare the result with the goal, and correct deviations. That is the part students often forget when they focus only on choice instead of execution.

\section*{6. UCM case template}
Use this sequence on any case. First, define the goal. Second, analyze the environment with PESTEL, Five Forces, and SWOT. Third, test the resources with RBV, VRIO, and the value chain. Fourth, check implementation and whether the structure actually supports the decision.

If the case is about growth, say whether the move is incremental or radical, what it changes, and what risks it creates. If the case is about defending a business, explain whether the firm should improve efficiency, reposition, or exit.

For the TechCore-style pivot case, the right questions are whether the firm should evolve, acquire, or pivot, and whether the move fits the market, the firm's resources, and the cost of execution.

The answer should compare at least two options, not just describe one. Then choose the one that best fits the firm's capabilities and the industry situation.

\section*{7. High-yield reminders}
If the industry looks unattractive but the firm has strong VRIO resources, the answer is usually some mix of differentiation, focus, or a protected niche.

If the industry is attractive but the firm lacks resources, then the right answer may be to build capabilities first, partner, or make a smaller entry rather than going all in.

If you recommend growth, explain why the move matches the firm's capabilities. If you recommend staying put, explain why the current position is defensible.

If you can, mention one downside of your own recommendation. That makes the answer sound more realistic and strategic.

The best answers from the slides tend to mention clusters, the value chain, the strategy clock, and the TechCore pivot because those examples show real application instead of memorized labels.

Clusters matter because geography, suppliers, institutions, and knowledge spillovers can improve productivity and innovation. That is a useful extra point when explaining why some firms outperform others.

Always end with one sentence on implementation risk. That keeps the answer strategic instead of purely analytical.

Implementation risk examples include weak leadership, resistance to change, insufficient resources, and poor KPI design.

\section*{Last-minute checklist}
Before you hand in the exam, check whether you have named the framework, used evidence, linked the evidence to the implication, and given a recommendation. If your answer is short on application, add one example from the case and one sentence on why the recommendation is feasible.

If you freeze, use this order: define the issue, pick the framework, write the strongest point first, then the second point, then the recommendation. Do not waste time listing definitions without interpretation.

\section*{Rapid revision table}
\begin{tabular}{@{}p{1.6cm}p{4.9cm}@{}}
\toprule
\term{Topic} & \term{Memory hook} \\
\midrule
Five Forces & Profit comes from the weakest force. \\
VRIO & Advantage needs value, rarity, inimitability, and organization. \\
Ansoff & New products and new markets raise risk quickly. \\
BCG & Invest in Stars, fund growth with Cash Cows, be careful with Question Marks. \\
BSC & What gets measured gets managed. \\
\bottomrule
\end{tabular}

\section*{Extra exam phrases}
The industry appears unattractive because the forces are strong and margins are pressured. The firm appears advantaged because its resources are valuable and hard to imitate. The strategy is coherent only if the environment, resources, and implementation all fit together. The best growth option is the one that uses existing capabilities while creating a realistic path to scale.

\end{multicols}
\end{document}